\chapter{Diseño Experimental}
\label{chap:diseno_experimental}

\section{Principios de diseño}
El diseño experimental se construyó bajo el principio de que cada experimento debe responder una pregunta concreta asociada a la hipótesis. En esta tesis se investigan cuatro preguntas: (i) si la regularización temporal mejora la reconstrucción bajo pérdida severa; (ii) si la ventaja se mantiene cuando los baselines clásicos son optimizados; (iii) si las mejoras se observan solo en error puntual o también en morfología y espectro; y (iv) qué limitaciones quedan abiertas al comparar implementaciones proxy con referencias externas.

\section{Mapa de hipótesis a experimentos}
\begin{longtable}{p{0.28\textwidth}p{0.30\textwidth}p{0.30\textwidth}}
\caption{Relación entre hipótesis operativas, experimentos y evidencia esperada.}\label{tab:hypothesis_map}\\
\toprule
Hipótesis operativa & Experimento & Evidencia esperada \\
\midrule
\endfirsthead
\toprule
Hipótesis operativa & Experimento & Evidencia esperada \\
\midrule
\endhead
TRSS mejora fidelidad temporal ante pérdida severa & Curvas de degradación por razón de canales faltantes & Menor DTW y mayor SNR en 30--40\% de pérdida \\
Los baselines optimizados reducen brechas aparentes & Optuna para splines, RBF, ICA y métodos GSP & Comparación más conservadora y auditable \\
La utilidad clínica requiere métricas espectrales & Análisis LSD, coherencia, PSD y bandas & Mejor preservación de estructura espectral \\
La pérdida contigua es más desafiante que la aleatoria & Escenarios \textit{nearby} vs \textit{random} & Mayor deterioro de métodos puramente espaciales \\
La equivalencia BGSRP estricta no está cerrada & Comparación cross-stack y reporte INS-13 & Limitación explícita con brecha residual cuantificada \\
La ventaja de TRSS debe sobrevivir a un baseline clínico-industrial fuerte & Benchmark confirmatorio TRSS vs. \code{MNE interpolate\_bads(method='spline')} & Mejora en error de amplitud sin ocultar coste computacional ni pérdidas en métricas espectrales \\
\bottomrule
\end{longtable}

\section{Datasets y justificación}
Los datasets se seleccionaron para cubrir distintas condiciones de densidad espacial y dinámica temporal. PhysioNet EEG Motor Movement/Imagery aporta una configuración de alta densidad con 64 canales; BCI Competition IV 2a aporta un régimen de menor densidad con 22 canales y mayor dificultad espacial; MNE Sample permite analizar eventos evocados y preservar morfologías ERP. Esta combinación evita que el benchmark quede restringido a un único montaje.

\section{Modos de pérdida de canales}
Se consideran dos familias de daño. En el modo \textit{random}, los canales faltantes se seleccionan de manera uniforme. Este modo modela fallas independientes o rechazo individual por artefactos. En el modo \textit{nearby}, se selecciona un canal semilla y luego sus vecinos espaciales más cercanos. Este modo modela fallas de contacto o desconexión local, que son más realistas cuando una región de la gorra EEG pierde calidad simultáneamente.

\section{Optimización y separación de datos}
La optimización de hiperparámetros se realiza con Optuna \cite{akiba2019optuna}. Para evitar fuga temporal, las señales se dividen en segmento de optimización y segmento de prueba. Entre ambos se introduce una zona de separación temporal cuando corresponde, reduciendo la posibilidad de que la fuerte autocorrelación del EEG haga artificialmente fácil la predicción del tramo de prueba. Los resultados declarados en esta tesis se interpretan como rendimiento bajo el protocolo de prueba definido, no como garantía universal fuera del espacio de búsqueda.

\section{Separación exploratoria--confirmatoria y benchmark contra MNE}

La tesis distingue entre una etapa exploratoria amplia y una etapa confirmatoria focalizada. La etapa exploratoria corresponde al pipeline multi-método B2/B3/B4 y a la optimización Optuna sobre familias geométricas, GSP instantáneas y espacio-temporales. Su objetivo fue mapear el espacio de métodos, detectar fallos de implementación, estimar rangos de hiperparámetros y generar hipótesis operativas. En cambio, el benchmark confirmatorio se diseñó después de observar que el baseline más relevante para uso práctico no es una spline reimplementada en forma aislada, sino \code{MNE interpolate\_bads(method='spline')}, por su estabilidad, documentación y adopción en análisis EEG.

El benchmark confirmatorio se ejecutó con pares comparables de condiciones y tres variantes de TRSS: configuración fija, configuración calibrada por validación cruzada con semilla cero y una variante oráculo de grilla usada solo como techo de rendimiento, no como estimador desplegable. La inferencia se reporta de forma pareada contra MNE, con 300 pares válidos, 100 clústeres de escenario definidos por dataset $\times$ modo de pérdida $\times$ nivel de pérdida, bootstrap jerárquico por clúster y corrección Benjamini--Hochberg sobre las métricas. Esta separación evita que las conclusiones finales dependan exclusivamente de iteraciones exploratorias o de comparaciones contra baselines sub-optimizados.

\section{Métricas y criterios de interpretación}
La evaluación usa métricas complementarias: MAE y RMSE para error de amplitud; SNR para calidad global; DTW para morfología temporal; LSD para distancia espectral logarítmica; y coherencia magnitud-cuadrada para consistencia de fase y frecuencia. La coexistencia de métricas evita seleccionar un método que minimice un error puntual destruyendo información fisiológica relevante.

\begin{table}[H]
\centering
\caption{Criterio de lectura de métricas empleadas en la tesis.}
\label{tab:metric_reading}
\begin{tabular}{p{0.16\textwidth}p{0.18\textwidth}p{0.45\textwidth}}
\toprule
Métrica & Dirección & Interpretación \\
\midrule
MAE & Menor mejor & Error absoluto medio de amplitud. \\
RMSE & Menor mejor & Penaliza errores grandes y picos espurios. \\
SNR & Mayor mejor & Relación señal-ruido reconstruida. \\
DTW & Menor mejor & Distancia morfológica robusta a desfases leves. \\
LSD & Menor mejor & Distancia entre espectros en escala logarítmica. \\
Coherencia & Mayor mejor & Consistencia frecuencial entre señal real y reconstruida. \\
\bottomrule
\end{tabular}
\end{table}

\section{Criterio estadístico}
Las comparaciones estadísticas se interpretan con pruebas no paramétricas cuando la distribución de métricas no puede asumirse normal. Se reportan valores $p$ y tamaños de efecto, privilegiando la lectura conjunta de significancia estadística y magnitud práctica. Un valor $p$ pequeño sin tamaño de efecto relevante no basta para justificar una conclusión fuerte; inversamente, una diferencia práctica puede requerir más datos si su incertidumbre es elevada.
